\section{Finite-index packets and the first exact cut}\label{sec:packets}

This section establishes the part of the finite-index geometry that does
not require the subsequent conditional-bijection groupoid.  All
relations below are actual invariant equivalences on the \emph{real}
sort of $M$.  Write $\Lambda$ for their finite lattice.  It is finite
because each member is a union of finitely many named raw atoms,
including equality.  We order it by inclusion; consequently $J\le K$
means that every $J$-class lies in a $K$-class.  For $J\le K$ let
$[K:J]$ denote the (constant, by transitivity) number of $J$-classes
contained in any $K$-class, possibly infinite.  Since all invariant
binary relations commute, the join of two invariant equivalences is
their represented product.

\begin{definition}[Finite-index closure and packet core]\label{def:packetcore}
For $J\in\Lambda$, let $\Sigma(J)$ be the greatest equivalence $K\ge J$
for which $[K:J]$ is finite.  Let $F=\Sigma(\Id)$.  An atomic diversity
colour $a$ is \emph{looped} if $(a,a,a)$ is an allowed triangle.  Write
$\Gamma(a)$ for the least invariant equivalence containing $a$, and put
\[
\kappa(J)=\bigvee\{\Gamma(a): a\subseteq J,\ (a,a,a)
                  \text{ is allowed}\},
\]
where an empty join is $\Id$.  Here and below $a\subseteq J$ means
that the entire actual $a$-orbital belongs to the raw support of $J$.
\end{definition}

Existence of $\Sigma(J)$ needs explanation: if $K_1,K_2\ge J$
have finite relative indices, their join $K_1K_2$ does too.  For each
fixed $K_1$-class, a $K_2$-class meeting it contains at most
$[K_2:J]$ distinct $J$-classes, while an alternating $K_1,K_2$-chain
can be rearranged into a single $K_1$-step followed by a single
$K_2$-step because the equivalences commute.  Thus
$[K_1K_2:J]\le[K_1:J][K_2:J]$.  Joining the finite family of all
finite-index coarsenings of $J$ gives $\Sigma(J)$.

\begin{lemma}[Loop criterion]\label{lem:loopcrit}
For invariant equivalences $J\le K$,
\[
 [K:J]<\infty\quad\Longleftrightarrow\quad
 \text{no looped diversity atom lies in }K\setminus J.
\]
\end{lemma}
\begin{proof}
If a looped $a$ lies in $K\setminus J$, then for every $n$ the
complete $a$-coloured graph on $n$ vertices belongs to $\Age(M)$:
all of its triangles are $(a,a,a)$.  Its vertices belong to one
$K$-class and to pairwise distinct $J$-classes.  Hence $[K:J]\ge n$
for every $n$.  Conversely, if one $K$-class contains infinitely
many $J$-classes, select an infinite family of representatives.
The infinite Ramsey theorem provides a monochromatic triangle on
three of them, in some colour $a$.  This colour is in $K\setminus J$
and is looped because the triangle occurs in $M$.
\end{proof}

\begin{lemma}[Canonical packet interval and exact core]
\label{lem:packetinterval}
For every $J\in\Lambda$,
\[
 [J:\kappa(J)]<\infty,
 \qquad \Sigma(\kappa(J))=\Sigma(J).
\]
The equivalences with finite-index closure $\Sigma(J)$ form precisely
the interval $[\kappa(J),\Sigma(J)]$; in particular $\kappa(J)$ is its
unique least element.  If $\kappa(J)>\Id$, then $\kappa(J)$ is an
actual positive word profile, and hence an exact equivalence.
\end{lemma}
\begin{proof}
Every looped colour of $J$ already belongs to $\kappa(J)$,
so the relative index is finite by Lemma~\ref{lem:loopcrit}.
If $J\le K$ has finite index, then the same lemma says that
$J$ and $K$ contain exactly the same looped colours.  Therefore
$\kappa(J)=\kappa(K)$, in particular for $K=\Sigma(J)$.
The finite-index closures of $J$ and $\kappa(J)$ consequently
coincide.  For any $L$ in their common packet one has
$\kappa(L)=\kappa(J)$ and $L\le\Sigma(L)=\Sigma(J)$.
Conversely, if $\kappa(J)\le L\le\Sigma(J)$, then
$[\Sigma(J):L]$ is finite, and therefore $\Sigma(L)=\Sigma(J)$.

To show exactness, fix a looped colour $a$.  Since the
$(a,a,a)$ triangle is allowed, $a\subseteq a^2$.  Symmetry also
gives $\Id\subseteq a^2$: every vertex has an $a$-neighbour by
transitivity and nonemptiness of the orbital.  The even powers
$a^{2m}$ therefore increase and, being unions of finitely many
actual orbitals, eventually stabilize to some relation $R$.
Associativity gives $R^2=R$; the relation is symmetric, contains
$\Id$ and $a$, and is therefore precisely $\Gamma(a)$.
Thus $\Gamma(a)$ is a positive word profile.  The finite join of
these equivalences is their product; if the resulting core is
nontrivial, at least one factor occurs and the product is again
a positive word profile.
\end{proof}

An invariant equivalence is called \emph{infinite} if its classes
are infinite.  If $M$ is infinite, the finite lattice $\Lambda$ has
an inclusion-minimal infinite member (possibly $\Top$).  Fix such a
member and denote it by $E$.  This is the \emph{minimal infinite cut}
used below.  In particular we do \emph{not} require $E<\Sigma(E)$:
that extra condition describes a different historical construction
(the minimal elevated packet core) and need not hold in general.

\begin{proposition}[Geometry below a minimal infinite cut]
\label{prop:minimalcut}
Let $E$ be an inclusion-minimal infinite invariant equivalence and
$H=E\cap F$.  Then $E=\kappa(E)$ is exact, and every proper
invariant equivalence $J<E$ has finite classes and lies in $H$.
In particular
\[
 [\Id,E]=[\Id,H]\cup\{E\},
\]
and $E$ admits no proper finite-index refinement.  If
$J<K\le E$ has finite index and $J<K$, then $K\le H$.
The actual quotient $M/E$ is a homogeneous triangle-defined AP and
semantic-homomorphic-CCO source in its binary orbital language.
\end{proposition}
\begin{proof}
Lemma~\ref{lem:packetinterval} gives $[E:\kappa(E)]<\infty$.
If the $\kappa(E)$-classes were finite, the $E$-classes would
also be finite, contradiction.  Hence $\kappa(E)$ is infinite;
minimality of $E$ yields $\kappa(E)=E$.  Exactness now follows
from Lemma~\ref{lem:packetinterval}.
For every $J<E$, minimality says that $J$ has finite classes.
By maximality of $F=\Sigma(\Id)$, this implies $J\le F$ and
therefore $J\le H$.  Conversely every $J\le H$ lies below $E$,
proving the displayed interval formula.  If $J<E$ had finite
index in $E$, its finite classes and the finite index would
make the classes of $E$ finite.  Thus no proper finite-index
refinement exists.  For a strict finite-index interval $J<K\le E$,
we must have $K<E$ and hence $K\le H$.  Since $E$ is exact,
the final quotient assertion is Theorem~\ref{thm:exact-sorts-source}.
\end{proof}

The next lemma handles intervals that cross the cut.  Its proof uses
only commutation of the invariant equivalences and minimality of $E$;
no assertion about the still-to-be-defined conditional groupoid is
being assumed.

\begin{lemma}[Relative finite-index identity]\label{lem:crossingindex}
Let $E$ be as in Proposition~\ref{prop:minimalcut} and put $H=E\cap F$.
If $J\le K$ and $[K:J]<\infty$, then
\begin{align}
 K\cap(E\vee J)&=K\cap(H\vee J),\label{eq:crossing-intersect}\\
 [K:J]&=[E\vee K:E\vee J]\,[E\cap K:E\cap J].
    \label{eq:crossing-index}
\end{align}
If $E\cap J<E\cap K$, then $E\cap K\le H$.
For every $J$ one also has
\begin{equation}\label{eq:sigma-crossing}
 \Sigma(J)\cap(E\vee J)=J\vee H.
\end{equation}
\end{lemma}
\begin{proof}
Write $I=E\cap J$ and $L=E\cap K$.  The map taking an
$I$-class in a fixed $L$-class to its containing $J$-class
shows $[L:I]\le[K:J]<\infty$.
If $I<L$, then $L\ne E$ by the absence of a proper
finite-index refinement of $E$.  Since every proper equivalence
below $E$ has finite classes, $L\le F\cap E=H$.

For commuting equivalences the modular identity
$K\cap(E\vee J)=J\vee(E\cap K)$ holds when $J\le K$:
a path of one $E$-step and one $J$-step connecting points
already in the same $K$-class can be rerouted inside that
$K$-class.  If $I=L$, then $L\le J$, so
$J\vee(E\cap K)=J$.  Moreover $K\cap H\le E\cap K=L$.
If $I<L$, then $L\le H$, so $K\cap H=L$.
In both cases
\[
 J\vee(E\cap K)=J\vee(K\cap H)=K\cap(J\vee H),
\]
proving~\eqref{eq:crossing-intersect}.

To count the classes, project the $J$-classes in one $K$-class
to the $(E\vee J)$-classes in its $(E\vee K)$-class.  This
projection is onto by permutability: a specified $(E\vee J)$-class
inside the $(E\vee K)$-class of the chosen $K$-class contains
an element of that $K$-class, by rearranging an $E$-then-$K$ path.
Two $J$-classes inside that $K$-class project to the same point
precisely when they are related by $K\cap(E\vee J)=J\vee L$.
Within one $(J\vee L)$-class, its $J$-classes are in bijection
with the $I$-classes in one $L$-class, using permutability
and $I=J\cap L$.  Thus each fibre has cardinality $[L:I]$.
This gives~\eqref{eq:crossing-index}, and all factors are finite.

Finally use $K=\Sigma(J)$.  To justify $F\le\Sigma(J)$,
observe that $F\vee J$ has finite index over $J$ (at most the
finite size of an $F$-class), so $F\vee J\le\Sigma(J)$.
Thus $H\le F\le\Sigma(J)$, 
\eqref{eq:crossing-intersect} becomes
$\Sigma(J)\cap(E\vee J)=\Sigma(J)\cap(H\vee J)=H\vee J$,
which is~\eqref{eq:sigma-crossing}.
\end{proof}

\begin{remark}[Scope]\label{rem:packet-scope}
The assertions above establish the packet bottom, the actual minimal infinite cut (which may equal $\Top$), and the lattice identities.  They \emph{do not} construct the
conditional-bijection groupoid, prove relative real closure, or show
that a mixed typed quotient is automatically a new transitive CCO
source.  Those are separate obligations of the forthcoming chapters.
\end{remark}

\section{Finite exact regularization and the thin bottom}
\label{sec:finite-regularization}
An atomic diversity relation $a$ is \emph{thin} if for every
$x\in M$ there is exactly one $y$ with $x\mathrel a y$.
By transitivity its valency is independent of $x$.

\begin{lemma}[Finite exact regularization]\label{lem:finite-regularization}
Let $F=\Sigma(\Id)$, the greatest invariant equivalence with
finite real classes.  There is an exact equivalence $H_0\le F$,
with finite classes, such that all atomic relations lying below
the induced equivalence $F/H_0$ on $M/H_0$ are thin.
We may take $H_0=\Id$ if all atoms within $F$ are thin.
If an atom within $F$ is nonthin, we may take $H_0>\Id$.
\end{lemma}
\begin{proof}
List the finitely many actual atoms, including equality, whose
supports lie in $F$: $a_0=\Id,a_1,\ldots,a_r$.
Every $a_i^2$ contains $\Id$, since $a_i$ is symmetric and has
a neighbour at every point.  Form the represented word profile
$P=a_1^2\cdots a_r^2$, with $P=\Id$ if the list is empty.
It is symmetric and contains $\Id$.  Consequently its powers
$P,P^2,\ldots$ increase and eventually stabilize in the finite
palette, say to $H_0=P^n=P^{n+1}$.
Then $H_0^2=H_0$, so $H_0$ is reflexive, symmetric and transitive.
Because $P\le F$ and $F$ is an equivalence, $H_0\le F$, hence its
classes are finite.  As a represented word profile, $H_0$ is exact.
Since $P$ contains each $a_i^2$, one has $a_i^2\le H_0$.

The $H_0$-saturated pullback $H_0a_iH_0=H_0a_i$ is one
actual atomic pair orbital of the exact quotient $M/H_0$
(Theorem~\ref{thm:joint-exact}).  Its relational square is
$H_0a_i^2H_0=H_0$ because $a_i^2\le H_0$ and $\Id\le a_i^2$.
Thus the induced quotient orbital has square equal to the diagonal.
If a nonempty symmetric binary orbital $b$ satisfies
$b^2=\Id$, then any two $b$-neighbours of one vertex are equal;
transitivity gives existence of a neighbour at every vertex.
Hence $b$ is thin.  Every atom inside $F/H_0$ arises in this
way from some $a_i$, proving the assertion.

If all $a_i$ are thin, then $a_i^2=\Id$ and $P=\Id$, so
$H_0=\Id$.  Conversely, if some $a_i$ is nonthin, two distinct
$a_i$-neighbours of a point give a non-diagonal pair in
$a_i^2\le H_0$, so $H_0>\Id$.
\end{proof}

\begin{lemma}[Geometric finite-exact-kernel normalisation]
\label{lem:finite-exact-normalization-geometric}
Let $M$ be the standing homogeneous symmetric triangle-defined
AP+CCO source in its finite binary orbital language.  There is a
finite chain of actual exact quotients
\[
              M=M_0\longrightarrow M_1\longrightarrow\cdots
                 \longrightarrow M_m
\]
(where $m$ may be zero) such that each nontrivial step is the
quotient of $M_i$ by a nonidentity invariant exact equivalence
$K_i$ with finite real classes.  Each $M_i$ is again a homogeneous
symmetric triangle-defined AP+CCO source.  The terminal source
$M_m$ has no invariant finite-class exact equivalence strictly
above $\Id$, and all its atomic orbitals contained in its
greatest finite-class equivalence $F_m=\Sigma_{M_m}(\Id)$ are thin.
\end{lemma}
\begin{proof}
If a nonidentity invariant finite-class exact equivalence $K$
exists in the current source, take its actual exact quotient.
By Theorem~\ref{thm:exact-sorts-source} the quotient retains all
standing source properties.  Since $K>\Id$, some original
diversity orbital becomes quotient equality, while each quotient
orbital is a union of original orbital supports.  Thus the number
of diversity orbitals strictly decreases, so only finitely many
such steps are possible.  At the terminal source no nontrivial
finite-class exact equivalence remains.  If an atomic orbital
contained in its greatest finite-class equivalence $F_m$ were
nonthin, Lemma~\ref{lem:finite-regularization} would produce a
nonidentity exact equivalence with finite classes, contradiction.
\end{proof}

\begin{example}[A thin finite bottom with a nontrivial exact kernel]
\label{ex:thin-not-exact-free}
The stronger geometric reduction in Lemma~\ref{lem:finite-exact-normalization-geometric} is really
necessary, even for a homogeneous triangle-defined AP+CCO source.
Let $I,J$ be countably infinite, and consider
$M=I\times J\times\{0,1\}$.  Name four diversity colours:
\begin{itemize}
\item $t$ between the two distinct points of one cell $(i,j)$;
\item $r$ if the $I$-coordinates agree but the $J$-coordinates differ;
\item $c$ if the $J$-coordinates agree but the $I$-coordinates differ;
\item $b$ if both coordinates differ.
\end{itemize}
The final two-point coordinate is immaterial for $r,c,b$.
The two equivalences given by equality of the $I$-coordinate and
of the $J$-coordinate are determined by the unions
$\Id\cup t\cup r$ and $\Id\cup t\cup c$ respectively.  Their
intersection is $F=\Id\cup t$, with classes of size two, and all
its atoms are thin.  The equivalences are exact: indeed $r^2$
and $c^2$ are precisely these two equivalences, so
$F=r^2\cap c^2$ is exact and nontrivial.

For completeness, the structure is homogeneous.  A finite partial
isomorphism induces partial bijections of the $I$- and $J$-index
sets, extendible to permutations of the respective countable sets;
within each occupied cell it induces a partial permutation of the
two-element fibre, independently extendible to a full permutation.
The age is triangle-defined: triangle constraints force the two
coordinate equality relations to be equivalences, while the forbidden
$t,t,t$ triangle enforces the bound of two points per cell.  Once
these conditions hold, distinct coordinate names and a two-point
labelling in each cell realise the complete finite diagram.  AP
follows from homogeneity.

The structure also has semantic homomorphic CCO.  For a partial
edge-coloured diagram, each specified edge requires equality or
inequality of the $I$-coordinates, and independently of the
$J$-coordinates, according to its colour.  Such a system of
constraints for either infinite coordinate fails precisely when a
specified inequality edge is joined by a path of specified equality
edges; a shortest such path plus the edge is a labelled simple cycle.
After these coordinate constraints are satisfied, the specified
$t$-edges are exactly the remaining requirements of different
$\{0,1\}$-labels inside one cell.  They can be satisfied precisely
when the $t$-edge graph is bipartite; an obstruction is an odd
monochromatic $t$-cycle.  Thus any failure of a homomorphic
completion is witnessed by a specified simple cycle, proving CCO.

Fix the infinite exact equivalence $E$ of equality of the
$I$-coordinate.  Between two distinct $E$-owners, the raw colour
$c$ has conditional valency $2$, because it chooses one $J$-coordinate
and either point of the target cell.  In particular
\[
                     E\cap c^2=F>\Id.
\]
This example refutes the shortcut ``$F$ thin, therefore there are no
nonidentity finite exact kernels'', but is not a counterexample to
EPPA.  The geometric reduction of Lemma~\ref{lem:finite-exact-normalization-geometric}
permits the first quotient by the exact kernel $F$ in this example;
the conditional-bijection argument then applies to the quotient.
\end{example}

\begin{lemma}[Elementary abelian translation group]
\label{lem:thin-bottom}
Assume every atom in $F$ is thin.  The atoms contained in $F$,
including equality, form a group $G\cong (\mathbb Z/2\mathbb Z)^d$
under represented relational composition.  For each $g\in G$
let $q_g$ be its atom and $t_g(x)$ the unique point with
$x\mathrel{q_g} t_g(x)$.  The $t_g$ are automorphisms of $M$,
commute with $\Aut(M)$ and act regularly on each $F$-class.
Moreover, for any original atomic colour $a$ the product $q_ga$
is a single atom, denoted $g\cdot a$, and
\begin{equation}\label{eq:thin-transport}
 \col(t_gx,t_hy)=(g+h)\cdot\col(x,y)
 \qquad(g,h\in G;\ x,y\in M),
\end{equation}
where equality is included among the possible colours.
\end{lemma}
\begin{proof}
Because every orbital below $F$ is a thin symmetric relation,
its unique-neighbour map $t_g$ is an involutive permutation.
For any $\sigma\in\Aut(M)$, uniqueness gives
$\sigma(t_gx)=t_g(\sigma x)$, so $t_g$ is equivariant.
The composite $t_gt_h$ is an equivariant permutation whose graph
is an invariant functional relation contained in $F$.  Such a
graph is a union of atomic orbitals.  Because each atom is
nonempty at every vertex, functionality forces precisely one
atom, and therefore $t_gt_h=t_k$ for a uniquely determined $k$.
This gives a finite group of commuting involutions with identity,
so $G$ is elementary abelian.  Every point of an $F$-class is
$t_gx$ for a unique $g$ by the atomic partition of that class.

The equivariant bijection $(x,y)\mapsto(t_gx,y)$ takes a single
$\Aut(M)$-orbital on ordered pairs bijectively to another one.
Thus $q_ga$ is exactly one atomic orbital, written $g\cdot a$.
Commutativity of invariant relational products gives the analogous
statement for switching the second endpoint and proves
\eqref{eq:thin-transport}.  Setting $h=g$ yields
$\col(t_gx,t_gy)=\col(x,y)$ since $g+g=0$,
so $t_g$ is an automorphism.  The equivariance already proved
shows that it is central in $\Aut(M)$.
\end{proof}

\begin{corollary}[Lower interval after thin regularization]
\label{cor:lower-subgroups}
In the reduced source in which $F$ is thin, let $E$ be an
inclusion-minimal infinite invariant equivalence and let
$H=E\cap F$.  Identifying the thin atoms in $F$ with the group
$G$ of Lemma~\ref{lem:thin-bottom}, the interval below $E$ is
\[
 [\Id,E]=\operatorname{Sub}(H)\cup\{E\}.
\]
Here $H$ is identified with a subgroup of $G$.
\end{corollary}
\begin{proof}
An invariant relation below $F$ is a union of the graphs of
translations $t_g$, and it is an equivalence exactly when its
selected subset of $G$ is a subgroup (closed under product and
inverse, containing the identity).  The equivalence $H\le F$
therefore corresponds to a subgroup, again denoted $H$.
Apply Proposition~\ref{prop:minimalcut}, which says that all proper
invariant equivalences below $E$ lie below $H$.
\end{proof}

\begin{remark}[Scope of the finite bottom]
\label{rem:finite-bottom-scope}
The regular action in Lemma~\ref{lem:thin-bottom} is a statement
about $F$-fibres in the \emph{actual} reduced homogeneous source.
A quotient by an arbitrary nonexact equivalence or a freely
selected system of packet origins need not inherit this action
in a compatible mixed-sort language.  The conditional groupoid
and the choice of all actual owners must be handled separately.
\end{remark}
